\documentclass[a4paper,20pt]{article}
\usepackage{latexsym}
\usepackage[empty]{fullpage}
\usepackage{titlesec}
\usepackage{marvosym}
\usepackage[usenames,dvipsnames]{color}
\usepackage{verbatim}
\usepackage{enumitem}
\usepackage{hyperref}
\usepackage{ctex}
\usepackage{fancyhdr}
\usepackage{needspace}

\pagestyle{fancy}
\fancyhf{}
\fancyfoot{}
\renewcommand{\headrulewidth}{0pt}
\renewcommand{\footrulewidth}{0pt}

% Adjust margins
\addtolength{\oddsidemargin}{-0.530in}
\addtolength{\evensidemargin}{-0.375in}
\addtolength{\textwidth}{1in}
\addtolength{\topmargin}{-.45in}
\addtolength{\textheight}{1in}

\urlstyle{rm}

\raggedbottom
\raggedright
\setlength{\tabcolsep}{0in}

% Sections formatting
\titleformat{\section}{
  \vspace{-10pt}\scshape\raggedright\large
}{}{0em}{}[\color{black}\titlerule \vspace{-6pt}]

%-------------------------
% Custom commands
\newcommand{\resumeItem}[2]{
  \item\small{
    \textbf{#1}{: #2 \vspace{-2pt}}
  }
}

\newcommand{\resumeItemWithoutTitle}[1]{
  \item\small{
    {\vspace{-2pt}}
  }
}

\newcommand{\resumeSubheading}[4]{
  \vspace{-1pt}\item
    \begin{tabular*}{0.97\textwidth}{l@{\extracolsep{\fill}}r}
      \textbf{#1} & #2 \\
      \textit{#3} & \textit{#4} \\
    \end{tabular*}\vspace{-3pt}
}

\newcommand{\resumeSubItem}[2]{\resumeItem{#1}{#2}\vspace{-3pt}}

\renewcommand{\labelitemii}{$\circ$}

\newcommand{\resumeSubHeadingListStart}{\begin{itemize}[leftmargin=*]}
\newcommand{\resumeSubHeadingListEnd}{\end{itemize}}
\newcommand{\resumeItemListStart}{\begin{itemize}}
\newcommand{\resumeItemListEnd}{\end{itemize}\vspace{-3pt}}

%-----------------------------
%%%%%%  CV STARTS HERE  %%%%%%

\begin{document}

%----------HEADING-----------------
\begin{tabular*}{\textwidth}{l@{\extracolsep{\fill}}r}
  \textbf{{\LARGE 黄业琦}} & 邮箱: \href{mailto:}{yeqi.huang@ed.ac.uk}\\
  \href{https://yeqi-huang.com}{个人主页: yeqi-huang.com}
\end{tabular*}

%-----------EDUCATION-----------------
\section{教育经历}
\resumeSubHeadingListStart
    \resumeSubheading{爱丁堡大学}{英国爱丁堡}
      {计算机科学博士}{2023年8月 - 至今}
\vspace{-3pt}
      {\scriptsize \textit{ \footnotesize{\newline{}\textbf{研究方向:} AI系统、分布式机器学习、面向ML的体系结构、无服务器计算}}}
\resumeSubheading{中国科学技术大学}{中国合肥}
      {计算机科学学士；少年班学院}{2017年7月 - 2021年6月}
\vspace{-3pt}
      {\scriptsize \textit{ \footnotesize{\newline{}\textbf{核心课程:} 操作系统、人工智能、编译原理、计算机体系结构、高性能体系结构}}}
\vspace{-3pt}
\resumeSubHeadingListEnd
\vspace{-10pt}

%-----------PUBLICATIONS-----------------
\section{学术论文}
\resumeSubHeadingListStart
\item [1] \textbf{Yeqi Huang}, Congjie He, Haocheng Xiao 等. ``Wavel: A Fast and Efficient Compilation System for Wafer-Scale Accelerators.'' \textit{已被ACM操作系统原理研讨会（SOSP）录用}, 2026.
\vspace{-5pt}
\item [2] Congjie He, Le Xu, Zhan Lu, \textbf{Yeqi Huang} 等. ``MeshRT: Compile-Time Governed Wafer-Scale Runtime for Low-Latency High-Throughput Inference.'' \textit{已被ACM操作系统原理研讨会（SOSP）录用}, 2026.
\vspace{-5pt}
\item [3] \textbf{Yeqi Huang}, Yanwei Ye, Guomin Chen 等. ``SwarmPilot: A Scheduler Agent Framework for Large Agentic Workflow Clusters.'' \textit{已被欧洲计算机系统会议（EuroSys）录用}, 2026.
\vspace{-5pt}
\item [4] \textbf{Yeqi Huang}, Yue Chen, Yanwei Ye 等. ``Ryze: Evidence-Enriched Data Synthesis from Biomedical Papers.'' 发表于 \textit{第64届计算语言学协会年会（ACL 2026）系统演示论文集}, 2026.
\vspace{-5pt}
\item [5] Chuanming Zha, Jiamin Zheng, \textbf{Yeqi Huang} 等. ``LLM-Monitor: Efficient Privacy Violation Monitoring for LLMs.''
\vspace{-5pt}
\item [6] Yinsicheng Jiang\textsuperscript{*}, \textbf{Yeqi Huang}\textsuperscript{*}, Liang Cheng 等. ``ContextPilot: Fast Long-Context Inference via Context Reuse.'' 发表于 \textit{机器学习与系统会议（MLSys）论文集}, 2026. (\textsuperscript{*}共同一作)
\vspace{-5pt}
\item [7] Congjie He, \textbf{Yeqi Huang}, Pei Mu 等. ``WaferLLM: Large Language Model Inference at Wafer Scale.'' 发表于 \textit{第19届 USENIX 操作系统设计与实现研讨会（OSDI 25）}, 2025.
\vspace{-5pt}
\item [8] Yinsicheng Jiang\textsuperscript{*}, Yao Fu\textsuperscript{*}, \textbf{Yeqi Huang}\textsuperscript{*} 等. ``MoE-CAP: Benchmarking Cost, Accuracy and Performance of Sparse Mixture-of-Experts Systems.'' 发表于 \textit{神经信息处理系统大会（NeurIPS）论文集}, 2025. (\textsuperscript{*}共同一作)
\vspace{-5pt}
\item [9] Yao Fu, Leyang Xue, \textbf{Yeqi Huang} 等. ``ServerlessLLM: Low-Latency Serverless Inference for Large Language Models.'' 发表于 \textit{第18届 USENIX 操作系统设计与实现研讨会（OSDI 24）}, 2024.
\vspace{-5pt}
\item [10] Xiao-Long Chen, Lin-Feng Wang, \textbf{Yeqi Huang} 等. ``Symplectic Structure-Preserving Particle-in-Cell Whole-Volume Simulation of Tokamak Plasmas.'' 发表于 \textit{SC21: 国际高性能计算、网络、存储与分析会议}, 2021.
\resumeSubHeadingListEnd



\section{项目经历}
\resumeSubHeadingListStart
\resumeSubItem{Wavel（SOSP 2026）}{设计MeshIR与局部性保持的晶圆级调度搜索，实现最高3.6倍加速，并将搜索时间从数周缩短至数秒。}
\resumeSubItem{MeshRT（SOSP 2026）}{将动态LLM推理事件编译为去中心化的逐核调度，在保持超低延迟的同时将吞吐量提升30--68倍。}
\resumeSubItem{SwarmPilot（EuroSys 2026）}{开发结合神经预测器、记忆、工具与协作机制的调度智能体；生产部署使P99延迟最高降低50\%，吞吐量翻倍。}
\resumeSubItem{Ryze（ACL 2026 Demo）}{构建证据增强的生物医学VLM数据流水线，并对Qwen3-VL-8B进行SFT与RL后训练，以低于200美元成本在LAB-Bench上超过GPT-5.2。}
\resumeSubItem{LLM-Monitor}{将隐私策略编译为高效的本地模型监控器；基于Qwen3-8B的专用模型以15倍速度达到GPT-5.2监控质量。}
\resumeSubItem{AutoReader}{构建科学文献的图与向量检索系统，索引速度比NVIDIA cuVS快60倍，检索速度比Milvus快13倍。}
\resumeSubItem{AI4Math - Sketchpad}{构建获AI for Math基金资助的系统，将数学证明转化为结构化图并分解，以支持更精确的自动形式化。}
\resumeSubItem{GPU上的QED模拟}{实现CUDA蒙特卡罗模拟，获得1500--5000倍加速，并在APS会议上获得诺贝尔奖得主Frank Wilczek认可。}

\resumeSubHeadingListEnd
\vspace{-10pt}

%----------------------------------------------------------------------------------------
%	TALKS
%----------------------------------------------------------------------------------------

\section{学术报告}
{\small
\resumeSubHeadingListStart
\item [1] \textbf{Yeqi Huang}. ``InfiniTensor: A Tensor-Friendly, Efficient Parallel Programming Library for Accelerator-Centric Clusters.'' \textit{UKSys}, 2024.
\vspace{-5pt}
\item [2] \textbf{Yeqi Huang}. ``Why we need a new clustering benchmark in AI retrieval?'' \textit{Int'l Workshop on Efficient Generative AI}, 2024.
\resumeSubHeadingListEnd
}

%----------------------------------------------------------------------------------------
%	TEACHING EXPERIENCE
%----------------------------------------------------------------------------------------

\section{教学经历}

\resumeSubHeadingListStart

\resumeSubItem{机器学习系统助教 (2025)}{
\begin{description}[font=$\circ$]\vspace{-8pt}
\item 担任CUDA编程课程助教，涵盖GPU架构、CUDA、CuPy和Triton。\vspace{-3pt}
\item 设计动手实验和项目，帮助学生掌握并行编程概念。\vspace{-3pt}
\end{description}
}
\resumeSubItem{机器学习系统助教 (2026)}{
\begin{description}[font=$\circ$]\vspace{-8pt}
\item 担任助教，新增CuTile编程模块，教授基于Tile的GPU编程。\vspace{-3pt}
\end{description}
}
\resumeSubHeadingListEnd

%----------------------------------------------------------------------------------------
%	AWARDS
%----------------------------------------------------------------------------------------
\section{获奖经历}
\textbf{国际奖项:}
\vspace{-8pt}
\begin{description}[font=$\bullet$]
\item {国际超级计算大会学生集群竞赛冠军 - 2021}
\vspace{-8pt}
\item {亚洲超算大会学生集群竞赛一等奖 - 2021}
\vspace{-8pt}
\end{description}
\textbf{国内奖项:}
\vspace{-8pt}
\begin{description}[font=$\bullet$]
\item {年度最佳中国超算应用 - 2022}
\vspace{-8pt}
\item {华为先锋开发者（全国仅4人）- 2021}
\vspace{-8pt}
\item {全国编译系统设计赛二等奖 - 2021}
\vspace{-8pt}
\end{description}

%----------------------------------------------------------------------------------------
%	SKILLS
%----------------------------------------------------------------------------------------

\section{专业技能}

\resumeSubHeadingListStart
\resumeSubItem{LLM系统与推理}{高性能服务、分布式与晶圆级推理、批处理与MoE动态运行时、模型加载与迁移、KV/上下文复用，以及GPU、Apple Silicon和Cerebras上的跨硬件优化。}
\resumeSubItem{LLM训练与评测}{证据增强数据合成、SFT与RL后训练、推理时扩展，以及面向LLM/VLM的质量、延迟、吞吐量、能效和成本评测。}
\resumeSubItem{RAG与智能体}{向量与图检索、上下文索引与去重、多智能体工作流、神经预测、记忆与工具抽象，以及Ray和ComfyUI生产集成。}
\resumeSubItem{系统与编程}{\textbf{C++、Python、Rust}；CUDA、CuPy、Triton、CuTile、MLIR/LLVM、OpenMP、Linux/eBPF、分布式系统、编译工具与性能工程。}
\resumeSubItem{物理与数学}{数值方法、线性代数、计算流体力学、分子动力学、量子力学与量子电动力学。}
\resumeSubHeadingListEnd



\section{特殊经历}

\resumeSubHeadingListStart
\resumeSubItem{开源爱好者}{参与贡献GNOME、LAMMPS和LLVM等知名项目，自2019年起使用GitHub展示开发历程和创意。}
\resumeSubItem{经营科学主题咖啡馆}{利用软件开发收入，在学校附近创办了Quantum Coffee——一个鼓励学生跨学科探索和讨论科学话题的协作空间。}
\resumeSubItem{UNICEF慈善项目}{自2019年起，将个人全部收入的10\%捐赠给儿童慈善机构。}
\resumeItemListEnd

\end{document}
